% Transcription — fonds Alexandre Grothendieck, shelfmark 14, batch 7 (pages 121-131, the last eleven pages of the folder).
% Established from the facsimile archives/batches/14/batch-07.pdf (a mirror of
% https://grothendieck.umontpellier.fr/14.pdf), per the skill
% transcribe-grothendieck. The batch file carries no cover sheet: PDF page N
% is page 120+N, as the Montpellier footer printed on each PDF page confirms,
% and the archivists' pencil numbers bottom-left agree with that position on
% every leaf (121 to 131).
%
% Pass: Opus 5.5 (claude-opus-5-5), 2026-09-23 — first pass, unchecked against the pages by a human.
%
% Transcribed: 122, 125, 127, 130, 131 (five pages). Skipped:
%   121, 123, 126, 128  leaves of SGA 3 typescripts used as paper — typed
%            pages 3 and 2 of an exposé on stabilisers (M_v(S'), the
%            diagram (D), « Lemme 1.1 ») and pages 24 bis and 25 of the
%            exposé on Cartan subgroups (« Corollaire 3.2 », « Corollaire
%            3.3 », « Lemme 3.4 », Chevalley), each cancelled by long
%            diagonal strokes. Their handwritten corrections concern the
%            typescript itself, not these notes; like the EGA typescript of
%            folder 12 batch 4, they are left out.
%   124, 129 blank leaves through which 125 and 128 show. Absent.
% No letter in this batch (the folder's 1965 letter is not among these pages).
%
% What the batch is. No heading and no pagination of his on any leaf; no date.
%   122      On the back of the typescript leaf 121: « Ceci prouve … les
%            implications (ii) ⇒ (i) ⇔ (iv) ⇒ (iii) ⇒ (v) », with three
%            glosses — the first « par construction … familles limitées de
%            faisceaux » (reading doubtful), (i) ⇔ (iv) by the definition of
%            K-equivalence of cycles, (iv) ⇒ (iii) by « R.R. ». It closes the
%            conditions (i)-(v) on a family M' ⊂ Z^i(X) of cycles (limited,
%            numerically limited, image of a limited family of locally free
%            sheaves via c^i, …) set out on page 120 (batch 6).
%   125      A chain of equivalent or successively implied statements on
%            cycles algebraically equivalent to zero: 1°) the product
%            A^i(X)^a × A^j(X)^a → A^{i+j}(X)^a is zero, 1° bis) the same
%            over X × Y; 2°), 2° bis) the same for zero-cycles A_0; 3°)
%            (x,y) − (x,b) − (a,y) + (a,b) ∼ 0, 3° bis) with X = Y; 4°) for
%            an abelian variety A, x ↦ (x) − (0) is a homomorphism A →
%            A_0(A) (« mais surjectif ! »); 4° bis) A_0(A)^a → A is an
%            isomorphism; 5°) A_0(X)^a → Alb is an isomorphism for every X;
%            6°), 6° bis) A_0(X×Y)^a → A_0(X)^a × A_0(Y)^a is bijective.
%   127      Scratch sheet for 125: the cycle (z × b) × (a × z') in Z × Z,
%            the map A → A_0(A) and its inverse, the scheme of implications
%            [2, 2 bis, 3, 3 bis] ⇔ [4, 4 bis] ⇒ [5], 5° ⇒ [6, 6 bis] ⇒ [2, …],
%            with a « th. de Weil » arrow and a « ? » converse both struck;
%            A_0(X×Y) → A_0(X) × A_0(Y) → Alb(X) × Alb(Y); φ(x) − φ(a).
%            Pencil jottings the other way up (group identities, a quotient
%            X²/R(S¹)) are given in a note.
%   130      A cycle of incidence Z ⊂ X × S over a rational variety S; C ⊂ X
%            of codim j+1, X' the blow-up with exceptional divisor P(N) ≃
%            C × S (as written), the fibre Y(y) over the generic point y, C_y
%            of codim j in it; x = f_*(y'), f^*f_*y' = y' + α_*(u), whence the
%            boxed A(Y(y)) = Im A(X) + Im A(C), and A^i(Y(y)) = Im A^i(X) +
%            Im A^{i−j}(C).
%   131      A few pencil jottings on the back of an IHÉS letterhead:
%            x² − y³, k[t], t², t³, k[x,y]/xy, and sketches of short
%            segments each carrying a point and of two crossing lines.
%
% Notation fixed for the batch. His script A of Chow groups is \mathcal{A},
% as in folder 12 and as his « Notations » of page 116 (batch 6) define it:
% A^*(X) = A^*_r(X) Chow ring, A_0 zero-cycles, the exponent a for classes
% algebraically equivalent to zero. On pages 125 and 127 some exponents are
% a letter like τ rather than his a; they are kept as τ and noted once. On
% page 130 the A of Chow groups is a plain capital and is transcribed so.
% His ʒ-shaped z on page 127 is z. His doubly underlined « nul » is \emph.
%
% Legibility. The formulas of 125 and 130 carry; the connective words of 122,
% 125 (4° bis, 5°) and the scratch of 127 are partly lost.
\documentclass[11pt,a4paper]{article}
% Le préambule commun à toutes les transcriptions du fonds Grothendieck.
%
% Il tient en une idée : **l'incertitude se compose, elle ne se résout pas.**
% Une transcription qui lisse un mot douteux a détruit la seule chose pour
% laquelle on la faisait — un lecteur doit pouvoir savoir, à chaque endroit,
% ce qui était sur le feuillet et ce qui a été ajouté. Les six macros qui
% suivent sont donc l'appareil critique, et non de la décoration.
%
% Les mêmes commandes sont reconnues par `scripts/render.mjs`, qui produit la
% vue de lecture du site. Une macro ajoutée ici doit l'être aussi là-bas,
% faute de quoi le rendu échouera bruyamment — ce qui est voulu : un rendu
% silencieusement faux est pire qu'un rendu absent.

\ProvidesPackage{grothendieck}[2026/08/08 Transcriptions du fonds Grothendieck]

\usepackage{amsmath,amssymb,amsthm}
\usepackage{mathrsfs}
\usepackage{stmaryrd}
% Les diagrammes commutatifs sont partout dans ce fonds : c'est la langue
% même de la géométrie algébrique de Grothendieck, et une transcription qui
% les rendrait en prose ne transcrirait rien.
\usepackage{tikz-cd}
% Français par défaut — c'est la langue des feuillets — mais l'anglais est
% chargé aussi : la traduction et le résumé sont des documents anglais, et la
% typographie française (espaces avant les deux-points) y serait une faute.
% Ces documents basculent par \selectlanguage{english} en tête de corps.
\usepackage[english,french]{babel}
\usepackage{fontspec}
\usepackage{geometry}
\geometry{a4paper,margin=25mm,marginparwidth=28mm}
\usepackage{marginnote}
\usepackage{xcolor}
% ulem, not soul. `soul`'s \st and \ul are built on a character-by-character
% scan that XeLaTeX's font machinery breaks: the document fails with an
% undefined control sequence rather than degrading. ulem does the same two jobs
% and is Unicode-engine safe. `normalem` stops it hijacking \emph, which the
% transcriptions use for Grothendieck's own emphasis.
\usepackage[normalem]{ulem}
\usepackage{hyperref}
\hypersetup{colorlinks=true,linkcolor=black,urlcolor=black,pdfborder={0 0 0}}

% --- Métadonnées du lot -----------------------------------------------------
% Reprises telles quelles par le script de rendu, qui les affiche en tête de
% la vue de lecture. Elles fixent ce que le fichier prétend couvrir : sans
% elles, une transcription flotte sans point d'attache dans un fonds de seize
% mille feuillets.

\newcommand{\folder}[1]{\gdef\grothFolder{#1}}
\newcommand{\batch}[1]{\gdef\grothBatch{#1}}
\newcommand{\leaves}[2]{\gdef\grothFirst{#1}\gdef\grothLast{#2}}
\newcommand{\foldertitle}[1]{\gdef\grothFoldertitle{#1}}
\gdef\grothFolder{?}\gdef\grothBatch{?}\gdef\grothFirst{?}\gdef\grothLast{?}\gdef\grothFoldertitle{}

% --- Le résumé --------------------------------------------------------------
%
% Il ouvre la lecture modernisée, sous le titre « Résumé », et il est composé
% en petit. Ce n'est pas un ornement : c'est le seul passage du document qui
% ne soit pas des mathématiques, et un lecteur doit pouvoir le reconnaître
% avant de l'avoir lu — puis le sauter s'il connaît déjà le sujet, ou s'y
% arrêter s'il ne le connaît pas. Le filet qui le suit marque la reprise du
% corps.

\newenvironment{resume}
  {\small\setlength{\parskip}{0.45em}}
  {\par\medskip\hrule\medskip}

% --- Mots-clés ---------------------------------------------------------------
%
% En anglais, en clôture du résumé de la lecture modernisée : le vocabulaire
% moderne sous lequel chercher ce que les feuillets construisent. Le manifeste
% du site (`npm run manifest`) les extrait pour étiqueter la cote entière —
% cette ligne est la source unique des tags, il n'y a pas de fichier à côté.
\newcommand{\keywords}[1]{\par\smallskip\noindent
  {\footnotesize\textsc{Keywords} --- \textit{#1}}\par}

% --- L'appareil critique ----------------------------------------------------

% Le numéro de feuillet, en marge. C'est l'ancre qui permet de régler un
% désaccord sur une formule en regardant *un* feuillet plutôt qu'une chemise
% de six cents. Le site s'en sert aussi pour tourner les pages du fac-similé
% quand on fait défiler la transcription.
\newcommand{\leaf}[1]{%
  \marginnote{\footnotesize\color{gray!70}\textsf{#1}}%
  \label{leaf:#1}\ignorespaces}

% Illisible : on ne devine pas. Un mot inventé qui se lit comme les autres est
% le pire résultat possible.
\newcommand{\ill}{\textcolor{red!60!black}{[\,\ldots\,]}}

% Lecture douteuse : le mot est donné, et signalé comme incertain.
\newcommand{\uncertain}[1]{\uline{#1}}

% Ajout de l'éditeur — une lettre manquante, un mot sous-entendu.
\newcommand{\add}[1]{\textcolor{blue!50!black}{[#1]}}

% Biffé par Grothendieck lui-même. Ce qu'il a rayé fait partie du document :
% une rature dit souvent où la pensée a changé de direction.
\newcommand{\struck}[1]{\sout{#1}}

% Note du transcripteur, sur l'état matériel du feuillet.
\newcommand{\note}[1]{\par\smallskip{\small\color{gray!60!black}\textit{[#1]}}\par\smallskip}

% Note marginale de Grothendieck, distincte du corps du texte.
\newcommand{\marginal}[1]{\par\smallskip\noindent\hspace{1em}%
  {\small\color{gray!60!black}#1}\par\smallskip}

% --- Titre ------------------------------------------------------------------

\AtBeginDocument{%
  \begin{center}
    {\large\bfseries Fonds Alexandre Grothendieck --- cote n\textsuperscript{o}~\grothFolder}\\[2pt]
    {\itshape\grothFoldertitle}\\[4pt]
    {\small feuillets \grothFirst--\grothLast\ (lot \grothBatch)}\\[2pt]
    {\footnotesize\color{gray!60!black}Transcription établie d'après le fac-similé numérisé
     par l'Université de Montpellier.\\
     Les lectures incertaines sont soulignées, les illisibles notées [\,\ldots\,],
     les ajouts entre crochets.}
  \end{center}
  \vspace{1em}}

\endinput


\folder{14}
\batch{7}
\pages{121}{131}
\dating{1965-[vers 1971]}
\watermark{Édition de démonstration}
\foldertitle{Théorie des cycles algébriques : conjectures de Hodge, Tate, Lefschetz, Weil : notes manuscrites (s.d.), lettre (1965).}

\begin{document}

\page{122}
\note{cette page conclut la liste des conditions (i) à (v) sur une famille $M'$ de cycles de $Z^{i}(X)$, posée à la page 120 ; la page 121, entre les deux, est un feuillet dactylographié étranger à ces notes}

Ceci prouve\supplied{,} \ill{}\add{,} les implications
\[
(\mathrm{ii}) \Longrightarrow (\mathrm{i}) \Longleftrightarrow (\mathrm{iv}) \Longrightarrow (\mathrm{iii}) \Longrightarrow (\mathrm{v})
\]
\marginal{par \uncertain{construction} \uncertain{que les} famille\supplied{s} \uncertain{limitées} de faisceaux}
\marginal{définition de \uncertain{la} K-équivalence de cycles}
\marginal{R.R.}
\note{les trois gloses sont écrites sous la chaîne d'implications, chacune rattachée par un trait à la flèche qu'elle justifie : la première à $(\mathrm{ii}) \Rightarrow (\mathrm{i})$, la deuxième à $(\mathrm{i}) \Leftrightarrow (\mathrm{iv})$, « R.R. » à $(\mathrm{iv}) \Rightarrow (\mathrm{iii})$ ; la première, en quatre lignes serrées, se lit mal. Le mot entre « prouve » et « les » est court et ne se lit pas}

\page{125}
\note{une accolade dans la marge de gauche embrasse les énoncés 2°) à 6° bis) ; entre deux énoncés consécutifs, une double flèche verticale, $\Updownarrow$ partout sauf entre 1° bis) et 2°), où elle est $\Downarrow$, et entre 5°) et 6°), où il n'y en a pas. Les mots « nul » sont soulignés deux fois}

1°) $\mathcal{A}^{i}_{\uncertain{r}}(X)^{a} \times \mathcal{A}^{j}(X)^{a} \longrightarrow \mathcal{A}^{i+j}(X)^{a}$ \quad \emph{nul}

$\Updownarrow$

1° bis) $\mathcal{A}^{i}(X)^{a} \times \mathcal{A}^{j}(Y)^{a} \longrightarrow \mathcal{A}^{i+j}(X \times Y)^{a}$ \quad \emph{nul}

$\Downarrow$

2°) $\mathcal{A}_{0}(X)^{a} \times \mathcal{A}_{0}(Y)^{a} \longrightarrow \mathcal{A}_{0}(X \times Y)^{a}$ \quad \emph{nul}

$\Updownarrow$

2° bis) $\mathcal{A}_{0}(X)^{a} \times \mathcal{A}_{0}(X)^{a} \longrightarrow \mathcal{A}_{0}(X \times X)^{a}$ \quad \emph{nul}

$\Updownarrow$

3°) $a, x \in X$, $b, y \in Y \Longrightarrow (x,y) - (x,b) - (a,y) + (a,b) \sim 0$
\note{sous « $a, x \in X$ » une petite insertion (un $a$ surligné, un signe $\in$ ?) ne se lit pas sûrement}

$\Updownarrow$

3° bis) Kif-kif avec $X = Y$.

$\Updownarrow$

\struck{$\mathcal{A}_{0}(X)^{\tau} \to \mathrm{Alb}(X)$}

4°) Si $A$ v.a.\ \ill{}, pour $x, y \in A$ :
\[
(x+y) - (x) - (y) + (0) \sim 0 ,
\]
i.e.\ $x \mapsto (x) - (0)$ est un \emph{homom.} $A \to \mathcal{A}_{0}(A)^{\tau}$ [mais surjectif !]
\note{le signe $\sim$ porte en dessous une petite marque, peut-être un $a$. L'exposant de $\mathcal{A}_{0}(A)$, ici comme dans la ligne biffée au-dessus de 4°) et à la page 127, a la forme d'un $\tau$ et non de son $a$ ; transcrit tel qu'il se lit}

$\Updownarrow$

4° bis) Si $A$ v.a., \struck{\ill{}} l'hom.
\[
\mathcal{A}_{0}(A)^{a} \longrightarrow A
\]
est un isom.\ (\uncertain{Donc} l'isom.\ réciproque \uncertain{est celui considéré dans} 4°)

$\Updownarrow$

5°) Pour tt $X$, \ill{} $\mathcal{A}_{0}(X)^{a} \longrightarrow \mathrm{Alb}(\uncertain{A})$ (qui est surj.) est un isomorphisme.
\note{la lettre dans $\mathrm{Alb}(\;)$ est repassée, un $A$ par-dessus un $X$ ou l'inverse}

6°) Pour $X, Y$, l'hom.\ $\mathcal{A}_{0}(X \times Y)^{a} \longrightarrow \mathcal{A}_{0}(X)^{a} \times \mathcal{A}_{0}(Y)^{a}$ (qui est surj.)\ est bijectif.

(6° bis) Kif-kif avec $X = Y$.

\page{127}
\note{feuille de brouillon pour la page 125, à l'encre et au crayon, sans ordre de lecture assuré ; transcrite de haut en bas}

\[
\begin{array}{cc} X & Y \\ a & b \end{array}
\qquad X \times Y = Z
\qquad\qquad A \times A \to A
\]

\[
(\underset{\sim\, 0}{z} \times b) \times (\underset{\sim\, 0}{a \times z'}) \quad \text{dans} \quad Z \times Z = X \times Y \times X \times Y
\]

\[
(z \times z') \times (b \times a)
\]

$3 \Longrightarrow 4$
\note{au crayon, à gauche, au-dessus de « $3 \Rightarrow 4$ », deux lignes écrites dans l'autre sens de la feuille, dont une biffée, se lisent peut-être $X^{2}/\mathbb{R}(S^{1}) \neq X^{2}(S^{1})/\mathbb{R}(S^{1})$ ; lecture très douteuse. Au crayon aussi et dans l'autre sens, à droite, une colonne d'identités : $a^{-1}(ab) = b$, $a(a^{-1}b) = b$, $(ba)a^{-1} = b$, $(ba^{-1})a = b$}

\struck{$\sum \mathcal{A}(x,y)$} \qquad $\sum u_{x,y}\,(x,y)$ \qquad \struck{$(x,y)$}

\begin{tikzcd}[column sep=huge]
A \arrow[r, "\text{hom. surjectif}"] & \mathcal{A}_{0}(A)^{\tau} \arrow[l, bend left=40, "\text{(hom. injectif)}"]
\end{tikzcd}

\note{le second mot de chacune des deux étiquettes est douteux (« surjectif », « injectif ») ; devant le $A$ de gauche, un astérisque repassé}

$A \longrightarrow \mathcal{A}_{0}(A)^{\tau}$ \quad \uncertain{est} bij.

\begin{tikzcd}
{[2, 2\,\mathrm{bis}, 3, 3\,\mathrm{bis}]} \arrow[r, leftrightarrow] & {[4, 4\,\mathrm{bis}]} \arrow[r, Rightarrow] & {[5]} \\
{[6, 6\,\mathrm{bis}]} \arrow[u, Rightarrow] & & \\
5^{\circ} \arrow[u, Rightarrow] & &
\end{tikzcd}

\marginal{\uncertain{via} les conditions}
\note{toutes les flèches de ce schéma sont doubles sur la page, celle entre les deux premiers crochets à double pointe. L'étiquette encadrée « via les conditions » est reliée par un trait à la flèche $\Leftrightarrow$ ; le « 5° » du bas est cerclé. Au-dessus de la flèche vers $[5]$, « \struck{th.\ de Weil} ». Une seconde double flèche montante, marquée « \struck{th.\ de Weil} », part d'un « 5° » biffé vers $[2, 2\,\mathrm{bis}, 3, 3\,\mathrm{bis}]$ ; une longue double flèche revient de la droite vers la gauche, marquée « ? » et accompagnée d'un texte biffé illisible. Deux grands traits au crayon entourent l'ensemble}

$\mathrm{Alb}(X \times Y)$

\struck{$\mathcal{A}_{0}(X \times Y) \to \mathcal{A}_{0}($}

$\mathcal{A}_{0}(X) \times \mathcal{A}_{0}(Y) \longrightarrow \mathcal{A}_{0}(X \times Y)$

\[
\mathcal{A}_{0}(X \times Y)^{\tau} \longrightarrow \mathcal{A}_{0}(X)^{\tau} \times \mathcal{A}_{0}(Y)^{\tau}
\]
$\downarrow$ \quad $\mathrm{Alb}(X) \times \mathrm{Alb}(Y)$
\note{la flèche verticale part de sous la flèche horizontale ; $\mathrm{Alb}(X) \times \mathrm{Alb}(Y)$ est écrit par-dessus un premier $\mathrm{Alb}(X \times Y)$, et la lecture de la surcharge est douteuse}

\struck{$x,y$} $x, y$ \qquad \struck{$\varphi(x) - \varphi(a)$}

\struck{$x,y$} \struck{$x - a$} \qquad $\varphi(x) - \varphi(a)$, $\varphi(y) - \varphi(b)$

$(x,y) - (x,b) - (a,y) + (a,b)$

$x - x - a + a$

\page{130}
$Z \subset X \times S$ \marginal{cycle d'incidence}

$Z^{i}(Y$
\note{l'expression s'arrête là}

$C \subset X \longleftarrow X'$, \quad $X' \to S$ \marginal{codim $j$} \marginal{éclaté}
\note{« codim $j$ » est écrit au-dessus de $C$, « éclaté » au-dessus de $X'$ ; un trait oblique relie ce schéma au $P(N)$ écrit plus haut à droite}

\begin{tikzcd}[column sep=large]
 & C \times S & \\
C \arrow[d, hook, "\operatorname{codim} j+1"'] & P(N) \arrow[l] \arrow[u, no head, "\simeq" description] \arrow[d, hook, "\alpha"] & C_{y} \arrow[l, hook'] \arrow[d, hook, "\alpha_{y}"] \\
X & X' \arrow[l, "f"] \arrow[d] & Y(y) \arrow[l, hook', "i'_{y}"] \arrow[d] \arrow[ll, bend left=35, "i_{y}"] \\
 & S & y \arrow[l]
\end{tikzcd}

\marginal{codim $j$} \marginal{pt générique} \marginal{variété rationnelle}
\note{« codim $j$ » est écrit à droite de $\alpha_{y}$, « pt générique » à droite de $y$, « variété rationnelle » sous $S$, qu'un trait lui rattache. $C \times S$ est écrit au-dessus de $P(N)$, relié par un signe $\simeq$ vertical ; la flèche courbe $i_{y}$ passe sous le diagramme, de $Y(y)$ à $X$. L'étiquette $f$ est placée sous la flèche $X \leftarrow X'$}

\[
y_{y} = i_{y}^{*}(y')
\]
\[
f_{*}(y') = x
\]
\[
i_{y}^{*}(x) = i'^{*}_{y}(f^{*}f_{*}(y')) \qquad f^{*}f_{*}y' = y' + \alpha_{*}(u)
\]
\[
\begin{aligned}
i_{y}^{*}(x) &= i'^{*}_{y}(y') + i'^{*}_{y}(\alpha_{*}(u)) \\ &= y' + \alpha_{y*}(u_{y})
\end{aligned}
\qquad \text{Or } u_{y} = u, \ u \in A(C)
\]
\note{à la dernière ligne il écrit $y'$ là où l'on attend $y_{y} = i'^{*}_{y}(y')$ ; transcrit tel quel}

Donc
\[
\boxed{A(Y(y)) = \mathrm{Im}\, A(X) + \mathrm{Im}\, A(C)}\,.
\]
\ill{}
\[
A^{i}(Y(y)) = \mathrm{Im}\, A^{i}(X) + \mathrm{Im}\, A^{i-j}(C)
\]
\note{le mot abrégé en tête de la dernière ligne ne se lit pas ; devant $A^{i-j}(C)$ un trait isolé, peut-être une parenthèse ébauchée}

\page{131}
\note{quelques notations au crayon, dans le sens de la hauteur}

\struck{$k[x$} \qquad $x^{2} - y^{\uncertain{3}}$ \qquad $k[t]$ \qquad $t^{2}, t^{3}$ \qquad $k[x,y]/xy$
\note{l'exposant de $y$ est repassé. À côté, des croquis : quatre courts segments portant chacun un point marqué, et deux droites qui se croisent}

\end{document}
